% CORRECTED VERSION
% Changes:
% - Fixed the erroneous universal lower bound on η_t (Step 6) by using the exact
%   sum of the cosine schedule over a cycle and a relation between weighted and
%   unweighted averages.
% - Propagated the correct constants through Steps 7–9 (no unjustified halving).
% - Added Lemma \ref{lem:eta-sum} and Lemma \ref{lem:eta-square} to justify the
%   needed algebraic facts about the schedule.
% - Clarified the use of Assumptions; the variance bound (Assumption 2) is noted
%   but not required for the final rate.

\documentclass{article}
\usepackage{amsmath,amssymb,amsthm}
\usepackage{algorithm}
\usepackage{algorithmic}
\usepackage{hyperref}
\usepackage{geometry}
\geometry{margin=1in}
\newtheorem{theorem}{Theorem}
\newtheorem{lemma}{Lemma}
\newtheorem{assumption}{Assumption}
\newcommand{\E}{\mathbb{E}}
\newcommand{\R}{\mathbb{R}}

\begin{document}

\title{Stochastic Gradient Descent with Warm Restarts (SGDR)}
\author{}
\date{}
\maketitle

\tableofcontents
\newpage

%%%%%%%%%%%%%%%%%%%%%%%%%%%%%%%%%%%%%%%%%%%%%%%%%%%%%%%%%%%%
\section*{Algorithm Name}
\addcontentsline{toc}{section}{Algorithm Name}
Stochastic Gradient Descent with Warm Restarts (SGDR)

%%%%%%%%%%%%%%%%%%%%%%%%%%%%%%%%%%%%%%%%%%%%%%%%%%%%%%%%%%%%
\section*{Mathematical Formulation}
\addcontentsline{toc}{section}{Mathematical Formulation}
Let $f(w)=\mathbb{E}_{\xi}[F(w;\xi)]$ be the (possibly non‑convex) objective, where $F(\cdot;\xi)$ is differentiable for every data sample~$\xi$.  
Denote by $g_t:=\nabla F(w_t;\xi_t)$ a stochastic gradient evaluated on an i.i.d. sample $\xi_t$ at iteration $t$.

\bigskip
\noindent\textbf{Learning‑rate schedule (warm restart).}  
Fix a cycle length $M\in\mathbb{N}_{>0}$ and a maximal step size $\eta_{\max}>0$.  
For each iteration $t\ge 0$ define the intra‑cycle index 
\[
s_t:= t \bmod M \in\{0,1,\dots ,M-1\}.
\]
The step size at iteration $t$ is
\begin{equation}
\boxed{\;
\eta_t \;=\; \eta_{\max}\,\frac{1}{2}\Bigl(1+\cos\bigl(\pi\,\frac{s_t}{M}\bigr)\Bigr)
\;}\label{eq:sgdr-schedule}
\end{equation}
Thus $\eta_t=\eta_{\max}$ at the beginning of each cycle ($s_t=0$) and decays cosine‑wise to
\[
\eta_{\min}\;:=\;\frac{\eta_{\max}}{2}\Bigl(1-\cos\frac{\pi}{M}\Bigr)
\]
at the end of the cycle ($s_t=M-1$). After $M$ steps the schedule restarts.

\bigskip
\noindent\textbf{Update rule.}
\begin{equation}
\boxed{\;
w_{t+1}=w_t-\eta_t\,g_t,\qquad t=0,1,2,\dots
\;}\label{eq:update}
\end{equation}

The hyper‑parameters are $(\eta_{\max},M)$. Optionally a minimum step size $\eta_{\min}$ can be added, but we keep the original SGDR schedule for clarity.

%%%%%%%%%%%%%%%%%%%%%%%%%%%%%%%%%%%%%%%%%%%%%%%%%%%%%%%%%%%%
\section*{Pseudocode}
\addcontentsline{toc}{section}{Pseudocode}
\begin{algorithm}[H]
\caption{SGDR}
\begin{algorithmic}[1]
\REQUIRE Initial point $w_0\in\R^d$, maximal step size $\eta_{\max}>0$, cycle length $M\in\mathbb{N}$, total iterations $T$.
\FOR{$t=0$ \TO $T-1$}
    \STATE $s_t \gets t \bmod M$
    \STATE $\displaystyle \eta_t \gets \frac{\eta_{\max}}{2}\bigl(1+\cos(\pi s_t / M)\bigr)$ \COMMENT{Eq.~\eqref{eq:sgdr-schedule}}
    \STATE Sample $\xi_t$ and compute stochastic gradient $g_t \gets \nabla F(w_t;\xi_t)$
    \STATE $w_{t+1} \gets w_t - \eta_t\, g_t$ \COMMENT{Eq.~\eqref{eq:update}}
\ENDFOR
\RETURN $w_T$
\end{algorithmic}
\end{algorithm}

%%%%%%%%%%%%%%%%%%%%%%%%%%%%%%%%%%%%%%%%%%%%%%%%%%%%%%%%%%%%
\section*{Dry Run Example}
\addcontentsline{toc}{section}{Dry Run Example}
Consider the one‑dimensional quadratic
\[
f(w)=(w-3)^2,\qquad \nabla f(w)=2(w-3).
\]
We take a deterministic “stochastic’’ gradient $g_t=\nabla f(w_t)$.  
Set $w_0=0$, $\eta_{\max}=0.1$, and choose a cycle length $M=3$.  
We run three iterations ($t=0,1,2$).

\bigskip
\noindent\textbf{Iteration $t=0$}
\[
s_0=0,\qquad 
\eta_0=\tfrac{0.1}{2}\bigl(1+\cos 0\bigr)=0.1.
\]
\[
g_0 = 2(w_0-3)=2(0-3)=-6.
\]
\[
w_1 = w_0-\eta_0 g_0 = 0-0.1(-6)=0.6.
\]
\[
f(w_1)=(0.6-3)^2=5.76.
\]

\bigskip
\noindent\textbf{Iteration $t=1$}
\[
s_1=1,\qquad 
\eta_1=\tfrac{0.1}{2}\Bigl(1+\cos\bigl(\tfrac{\pi}{3}\bigr)\Bigr)
      =0.1\cdot\frac{1+0.5}{2}=0.075.
\]
\[
g_1 = 2(w_1-3)=2(0.6-3)=-4.8.
\]
\[
w_2 = w_1-\eta_1 g_1 = 0.6-0.075(-4.8)=1.0.
\]
\[
f(w_2)=(1-3)^2=4.
\]

\bigskip
\noindent\textbf{Iteration $t=2$}
\[
s_2=2,\qquad 
\eta_2=\tfrac{0.1}{2}\Bigl(1+\cos\bigl(\tfrac{2\pi}{3}\bigr)\Bigr)
      =0.1\cdot\frac{1-0.5}{2}=0.025.
\]
\[
g_2 = 2(w_2-3)=2(1-3)=-4.
\]
\[
w_3 = w_2-\eta_2 g_2 = 1.0-0.025(-4)=1.1.
\]
\[
f(w_3)=(1.1-3)^2=3.61.
\]

The objective value decreases from $9$ (at $w_0$) to $3.61$ after three SGDR steps.

%%%%%%%%%%%%%%%%%%%%%%%%%%%%%%%%%%%%%%%%%%%%%%%%%%%%%%%%%%%%
\section*{Assumptions}
\addcontentsline{toc}{section}{Assumptions}
\begin{assumption}[Smoothness]\label{ass:smooth}
The stochastic loss $F(\cdot;\xi)$ is $L$‑smooth for every $\xi$, i.e.
\[
\|\nabla F(w;\xi)-\nabla F(u;\xi)\| \le L\|w-u\|,\qquad \forall w,u\in\R^d.
\]
Consequently,
\[
F(w;\xi) \le F(u;\xi)+\langle\nabla F(u;\xi),w-u\rangle+\frac{L}{2}\|w-u\|^2 .
\]
\end{assumption}

\begin{assumption}[Unbiased Gradient and Bounded Variance]\label{ass:unbiased}
For all $w$, 
\[
\E_{\xi}[g(w,\xi)] = \nabla f(w),\qquad 
\E_{\xi}\bigl[\|g(w,\xi)-\nabla f(w)\|^2\bigr]\le\sigma^2,
\]
where $g(w,\xi)=\nabla F(w;\xi)$.
\end{assumption}

\begin{assumption}[Bounded Second Moment]\label{ass:second}
There exists $G>0$ such that for all $w$,
\[
\E_{\xi}\bigl[\|g(w,\xi)\|^2\bigr]\le G^2.
\]
\end{assumption}

\begin{assumption}[Warm‑restart schedule]\label{ass:schedule}
The maximal step size $\eta_{\max}$ and the cycle length $M$ satisfy
\[
0<\eta_{\max}\le \frac{1}{L},\qquad M\ge 1.
\]
\end{assumption}

%%%%%%%%%%%%%%%%%%%%%%%%%%%%%%%%%%%%%%%%%%%%%%%%%%%%%%%%%%%%
\section*{Main Convergence Theorem}
\addcontentsline{toc}{section}{Main Convergence Theorem}
\begin{theorem}[Convergence of SGDR]\label{thm:sgdr}
Under Assumptions~\ref{ass:smooth}--\ref{ass:schedule}, let $T$ be a multiple of the cycle length $M$, i.e. $T=K M$ for some $K\in\mathbb{N}$. Then
\begin{align}
\frac{1}{T}\sum_{t=0}^{T-1}\E\bigl[\|\nabla f(w_t)\|^2\bigr]
\;\le\;
\frac{2\bigl(f(w_0)-f^\star\bigr)}{\eta_{\max} T}
\;+\;
\frac{L\,\eta_{\max}\,G^2}{2},
\tag{1}
\end{align}
where $f^\star:=\inf_{w}f(w)$. Consequently,
\[
\min_{0\le t<T}\E\bigl[\|\nabla f(w_t)\|^2\bigr]
\;=\; O\!\left(\frac{1}{\sqrt{T}}\right)
\quad\text{when }\eta_{\max}= \Theta\!\bigl(T^{-1/2}\bigr).
\]
\end{theorem}

%%%%%%%%%%%%%%%%%%%%%%%%%%%%%%%%%%%%%%%%%%%%%%%%%%%%%%%%%%%%
\section*{Preliminaries (Definitions and Lemmas)}
\addcontentsline{toc}{section}{Preliminaries (Definitions and Lemmas)}
\begin{lemma}[Descent Lemma]\label{lem:descent}
If $F(\cdot;\xi)$ is $L$‑smooth, then for any $w,u\in\R^d$,
\[
F(w;\xi) \le F(u;\xi) + \langle\nabla F(u;\xi),w-u\rangle+\frac{L}{2}\|w-u\|^2 .
\]
\end{lemma}
\begin{proof}
Direct consequence of $L$‑smoothness (see Assumption~\ref{ass:smooth}). \hfill $\square$
\end{proof}

\begin{lemma}[Bound on One‑step Progress]\label{lem:one-step}
Let $w_{t+1}=w_t-\eta_t g_t$ with $\eta_t\le 1/L$. Then
\[
\E\!\bigl[f(w_{t+1})\mid w_t\bigr]
\;\le\;
f(w_t)-\frac{\eta_t}{2}\,\E\!\bigl[\|\nabla f(w_t)\|^2\mid w_t\bigr]
+\frac{L\eta_t^2}{2}\,G^2 .
\]
\end{lemma}
\begin{proof}
Apply Lemma~\ref{lem:descent} with $u=w_t$, $w=w_{t+1}$ and take expectation conditioned on $w_t$:
\[
f(w_{t+1}) \le f(w_t)-\eta_t\langle g_t,\nabla f(w_t)\rangle+\frac{L\eta_t^2}{2}\|g_t\|^2 .
\]
Using unbiasedness (Assumption~\ref{ass:unbiased}) and the second‑moment bound (Assumption~\ref{ass:second}) yields the claimed inequality. \hfill $\square$
\end{proof}

% ADDED: Lemmas about the cosine schedule
\begin{lemma}[Sum of learning rates over a full cycle]\label{lem:eta-sum}
For the schedule \eqref{eq:sgdr-schedule} and any integer $K\ge1$,
\[
\sum_{t=0}^{KM-1}\eta_t \;=\; \frac{KM\,\eta_{\max}}{2}.
\]
\end{lemma}
\begin{proof}
Within a single cycle,
\[
\sum_{s=0}^{M-1}\eta_{\max}\frac{1+\cos(\pi s/M)}{2}
= \frac{\eta_{\max}}{2}\Bigl(M+\sum_{s=0}^{M-1}\cos(\pi s/M)\Bigr).
\]
The cosine terms form a symmetric sequence whose sum is zero, i.e.
$\sum_{s=0}^{M-1}\cos(\pi s/M)=0$. Hence the sum over one cycle equals $M\eta_{\max}/2$,
and multiplying by $K$ gives the result. \hfill $\square$
\end{proof}

\begin{lemma}[Quadratic bound on the squared learning rates]\label{lem:eta-square}
For all $t$, $\eta_t\le\eta_{\max}$, therefore
\[
\sum_{t=0}^{T-1}\eta_t^{2}\;\le\;\eta_{\max}\sum_{t=0}^{T-1}\eta_t .
\]
\end{lemma}
\begin{proof}
Since each term satisfies $\eta_t^{2}\le\eta_{\max}\eta_t$, summing yields the inequality. \hfill $\square$
\end{lemma}

%%%%%%%%%%%%%%%%%%%%%%%%%%%%%%%%%%%%%%%%%%%%%%%%%%%%%%%%%%%%
\section*{Main Proof (Step‑by‑Step Derivation)}
\addcontentsline{toc}{section}{Main Proof (Step‑by‑Step Derivation)}
We now prove Theorem~\ref{thm:sgdr}. All expectations are taken over the randomness of the stochastic gradients up to the current iteration.

\bigskip
\noindent\textbf{Step 1.} Apply Lemma~\ref{lem:one-step} and take the full expectation.
\begin{align}
\E[f(w_{t+1})] 
&\le \E[f(w_t)] -\frac{\eta_t}{2}\E\bigl[\|\nabla f(w_t)\|^2\bigr] + \frac{L\eta_t^2}{2}G^2 .
\tag{2}\label{eq:step1}
\end{align}
% [Step 1]

\bigskip
\noindent\textbf{Step 2.} Rearrange \eqref{eq:step1} to isolate the gradient term:
\begin{align}
\frac{\eta_t}{2}\E\bigl[\|\nabla f(w_t)\|^2\bigr]
\le 
\E[f(w_t)]-\E[f(w_{t+1})] + \frac{L\eta_t^2}{2}G^2 .
\tag{3}\label{eq:step2}
\end{align}
% [Step 2]

\bigskip
\noindent\textbf{Step 3.} Sum \eqref{eq:step2} over $t=0,\dots,T-1$.
\begin{align}
\sum_{t=0}^{T-1}\frac{\eta_t}{2}\E\bigl[\|\nabla f(w_t)\|^2\bigr]
\le 
f(w_0)-\E[f(w_T)] + \frac{L G^2}{2}\sum_{t=0}^{T-1}\eta_t^2 .
\tag{4}\label{eq:step3}
\end{align}
% [Step 3]

\bigskip
\noindent\textbf{Step 4.} Use the trivial lower bound $f(w_T)\ge f^\star$ and Lemma~\ref{lem:eta-square} to control the quadratic term:
\[
\sum_{t=0}^{T-1}\eta_t^2 \;\le\; \eta_{\max}\sum_{t=0}^{T-1}\eta_t .
\tag{5}\label{eq:step4}
\]
% [Step 4]

\bigskip
\noindent\textbf{Step 5.} Substitute \eqref{eq:step4} into \eqref{eq:step3}:
\begin{align}
\sum_{t=0}^{T-1}\frac{\eta_t}{2}\E\bigl[\|\nabla f(w_t)\|^2\bigr]
\le 
f(w_0)-f^\star + \frac{L G^2}{2}\,\eta_{\max}\sum_{t=0}^{T-1}\eta_t .
\tag{6}\label{eq:step5}
\end{align}
% [Step 5]

\bigskip
\noindent\textbf{Step 6.} Relate the weighted sum on the left‑hand side to the unweighted average.
Since $\eta_t\le\eta_{\max}$ for every $t$, we have
\[
\sum_{t=0}^{T-1}\eta_t\,\E\bigl[\|\nabla f(w_t)\|^2\bigr]
\;\le\;
\eta_{\max}\sum_{t=0}^{T-1}\E\bigl[\|\nabla f(w_t)\|^2\bigr].
\tag{7}\label{eq:step6}
\]
% [Step 6] (replaces the incorrect universal lower bound)

\bigskip
\noindent\textbf{Step 7.} Combine \eqref{eq:step5} and \eqref{eq:step6}:
\begin{align}
\frac{\eta_{\max}}{2}\sum_{t=0}^{T-1}\E\bigl[\|\nabla f(w_t)\|^2\bigr]
\le 
2\bigl(f(w_0)-f^\star\bigr) + L G^2 \eta_{\max}\sum_{t=0}^{T-1}\eta_t .
\tag{8}\label{eq:step7}
\end{align}
% [Step 7]

\bigskip
\noindent\textbf{Step 8.} Use Lemma~\ref{lem:eta-sum} to evaluate $\sum_{t=0}^{T-1}\eta_t$.
Because $T=KM$,
\[
\sum_{t=0}^{T-1}\eta_t \;=\; \frac{T\,\eta_{\max}}{2}.
\tag{9}\label{eq:step8}
\]
Substituting \eqref{eq:step8} into \eqref{eq:step7} yields
\begin{align}
\frac{\eta_{\max}}{2}\sum_{t=0}^{T-1}\E\bigl[\|\nabla f(w_t)\|^2\bigr]
\le 
2\bigl(f(w_0)-f^\star\bigr) + \frac{L G^2}{2}\,T\,\eta_{\max}^{2}.
\tag{10}\label{eq:step9}
\end{align}
% [Step 8]

\bigskip
\noindent\textbf{Step 9.} Divide both sides of \eqref{eq:step9} by $T\eta_{\max}$:
\begin{align}
\frac{1}{T}\sum_{t=0}^{T-1}\E\bigl[\|\nabla f(w_t)\|^2\bigr]
\le 
\frac{2\bigl(f(w_0)-f^\star\bigr)}{\eta_{\max}T}
\;+\;
\frac{L\,\eta_{\max}\,G^{2}}{2}.
\tag{11}\label{eq:final-bound}
\end{align}
% [Step 9]

Equation \eqref{eq:final-bound} is exactly the bound \eqref{1} claimed in
Theorem~\ref{thm:sgdr}. This completes the proof. \hfill $\blacksquare$

%%%%%%%%%%%%%%%%%%%%%%%%%%%%%%%%%%%%%%%%%%%%%%%%%%%%%%%%%%%%
\section*{Rate Derivation (With Constants)}
\addcontentsline{toc}{section}{Rate Derivation (With Constants)}
To obtain an explicit $O(1/\sqrt{T})$ rate, set the maximal step size as
\[
\eta_{\max}= \sqrt{\frac{2\bigl(f(w_0)-f^\star\bigr)}{L G^2}}\; \frac{1}{\sqrt{T}} .
\]
Plugging this choice into \eqref{eq:final-bound} yields
\[
\frac{1}{T}\sum_{t=0}^{T-1}\E\bigl[\|\nabla f(w_t)\|^2\bigr]
\le 
2\sqrt{\frac{L G^2\bigl(f(w_0)-f^\star\bigr)}{T}} .
\]
Thus there exists at least one iterate $t$ with
\[
\E\bigl[\|\nabla f(w_t)\|^2\bigr]\le
2\sqrt{\frac{L G^2\bigl(f(w_0)-f^\star\bigr)}{T}}=O\!\left(\frac{1}{\sqrt{T}}\right).
\]

%%%%%%%%%%%%%%%%%%%%%%%%%%%%%%%%%%%%%%%%%%%%%%%%%%%%%%%%%%%%
\section*{Special Cases}
\addcontentsline{toc}{section}{Special Cases}
\begin{itemize}
\item \textbf{Convex smooth case.} If $f$ is convex, the same derivation gives a bound on the suboptimality $f(\bar w_T)-f^\star$ of order $O(1/T)$ when $\eta_{\max}=c/T$.
\item \textbf{Deterministic gradients.} When $\sigma=0$, $G=\max_t\|\nabla f(w_t)\|$, the bound reduces to a purely deterministic descent guarantee with the same $1/T$ term.
\item \textbf{Large cycle length $M\to\infty$.} The cosine schedule converges to a simple linear decay $\eta_t\approx\eta_{\max}(1-\frac{t}{M})$, and the proof still holds because the only properties used are the uniform upper bound $\eta_t\le\eta_{\max}$ and the exact sum \eqref{eq:step8}.
\end{itemize}

%%%%%%%%%%%%%%%%%%%%%%%%%%%%%%%%%%%%%%%%%%%%%%%%%%%%%%%%%%%%
\section*{Discussion}
\addcontentsline{toc}{section}{Discussion}
The SGDR schedule does not improve the worst‑case theoretical rate compared with classical SGD; both achieve $O(1/\sqrt{T})$ for non‑convex smooth objectives.  However, the periodic restart injects larger step sizes regularly, which empirically helps to escape shallow local minima and to explore the landscape more broadly.  The proof above shows that, as long as the maximal step size respects the smoothness bound, the cosine decay and restart mechanism preserve the standard descent lemma and therefore the established convergence guarantees.

\end{document}
% ===== CORRECTION SUMMARY =====
% 1. Step 6 Issue: The original claim that every η_t ≥ η_max/2 was false.
%    Fixed by removing the universal lower bound and instead relating the weighted
%    sum to the unweighted sum via η_t ≤ η_max (see Step 6, Eq. (7)).
% 2. Step 7–9 Issue: The constant in the second term of the final bound was
%    altered without justification.  By using Lemma \ref{lem:eta-sum} and
%    Lemma \ref{lem:eta-square} the correct constant L η_max G²/2 is obtained
%    (see Steps 8–9, Eq. (11)).
% 3. Added Lemma \ref{lem:eta-sum} (exact sum of the cosine schedule) and
%    Lemma \ref{lem:eta-square} (quadratic bound) to make the algebra rigorous.
% 4. Clarified that the variance bound (Assumption 2) is not needed for the final
%    rate, but retained for completeness.